\documentclass[a4paper]{arrowhead}

\usepackage[yyyymmdd]{datetime}
\usepackage{etoolbox}
\usepackage[utf8]{inputenc}
\usepackage[T1]{fontenc}
\usepackage{multirow}
\usepackage{longtable}
\usepackage{tikz}
\usepackage{adjustbox}
\AtBeginEnvironment{longtable}{\small}
\BeforeBeginEnvironment{tikzpicture}{\begin{adjustbox}{max width=\textwidth}}
\AfterEndEnvironment{tikzpicture}{\end{adjustbox}}
\usetikzlibrary{positioning,shapes.multipart,arrows.meta}

\renewcommand{\dateseparator}{-}
\setlength{\emergencystretch}{3em}
\lstset{language=HTTP,breakatwhitespace=false,columns=fullflexible}

\newcommand{\fparam}[1]{\textit{\textcolor{ArrowheadBlue}{#1}}}

%% Special references
\newcommand{\mref}[2][]{{\textcolor{ArrowheadPurple}{\hyperref[sec:model:#2]{\ifstrempty{#1}{#2}{#1}}}}}
\newcommand{\oref}[1]{{\textcolor{ArrowheadBlue}{\hyperref[sec:operations:#1]{#1}}}}
\newcommand{\pdef}[1]{{\textcolor{ArrowheadGrey}{#1 \label{sec:model:primitives:#1} \label{sec:model:primitives:#1s}}}}
\newcommand{\pref}[1]{{\textcolor{ArrowheadGrey}{\hyperref[sec:model:primitives:#1]{#1}}}}

% as template_idd.tex; the input and output arguments are passed already formatted (\mref or \pref)
\newrobustcmd\osubsection[5]{
  \addtocounter{subsection}{1}
  \addcontentsline{toc}{subsection}{\protect\numberline{\thesubsection}operation \textcolor{ArrowheadBlue}{#1}}
  \subsection*{
    \thesubsection\quad
    #2 \textcolor{ArrowheadPurple}{#3} \\
    \small
    \hspace*{0.075\textwidth}\begin{minipage}{0.1\textwidth}
      \vspace*{1mm}
      Operation: \\
      \notblank{#4}{Input: \\}{}
      \notblank{#5}{Output: \\}{}
    \end{minipage}%
    \begin{minipage}{0.825\textwidth}
      \vspace*{1mm}
      \textcolor{ArrowheadBlue}{#1} \\
      \notblank{#4}{#4 \\}{}
      \notblank{#5}{#5 \\}{}
    \end{minipage}
  }
  \label{sec:operations:#1}
}
\newrobustcmd\msubsection[2]{
  \addtocounter{subsection}{1}
  \addcontentsline{toc}{subsection}{\protect\numberline{\thesubsection}#1 \textcolor{ArrowheadPurple}{#2}}
  \subsection*{\thesubsection\quad#1 \textcolor{ArrowheadPurple}{#2}}
  \label{sec:model:#2}
}
%%

\begin{document}
\sloppy

%% Arrowhead Document Properties
\ArrowheadTitle{authorizationToken MQTT/TLS/JSON}
\ArrowheadServiceID{authorizationToken}
\ArrowheadType{Interface Design Description}
\ArrowheadTypeShort{IDD}
\ArrowheadVersion{5.2.2}
\ArrowheadDate{2026-10-22}
\ArrowheadAuthor{Jerker Delsing}
\ArrowheadStatus{RELEASE CANDIDATE}
\ArrowheadContact{jerker.delsing@ltu.se}
\ArrowheadFooter{\href{www.arrowhead.eu}{www.arrowhead.eu}}
\ArrowheadSetup
%%

%% Front Page
\begin{center}
  \vspace*{1cm}
  \huge{\arrowtitle}

  \vspace*{0.2cm}
  \LARGE{\arrowtype}
  \vspace*{1cm}

  \Large{Service ID: \textit{"\arrowid"}}

  \vspace*{\fill}

  % Front Page Image

  \vspace*{1cm}
  \vspace*{\fill}

  % Front Page Abstract
  \begin{abstract}
    This document describes the MQTT protocol with TLS payload security and JSON payload encoding variant of the \texttt{authorizationToken} service of the ConsumerAuthorization system.
  \end{abstract}

  \vspace*{1cm}
\end{center}

\newpage
%%

%% Table of Contents
\tableofcontents
\newpage
%%

\section{Overview}
\label{sec:overview}

This document describes the \texttt{authorizationToken} service interface of the ConsumerAuthorization system over MQTT, which provides authorization tokens: a consumer obtains a token for a specific provider, service and operation; the provider verifies it.

It's implemented using protocol, encoding as stated in below table:
\begin{table}[ht!]
  \centering
  \begin{tabular}{|l|l|l|}
    \hline
    \rowcolor{gray!33} Profile type & Type & Version \\ \hline
    Transfer protocol & MQTT & 3.1.1 \cite{mqtt311} \\ \hline
    Data encryption & TLS & 1.3 \cite{rfc8446}, 1.2 \cite{rfc5246} \\ \hline
    Encoding & JSON & RFC 8259 \cite{rfc8259} \\ \hline
    Compression & N/A & - \\ \hline
    Semantics & N/A & - \\ \hline
    Ontology & N/A & - \\ \hline
  \end{tabular}
  \caption{Communication and semantics details used for the \texttt{authorizationToken} service interface}
  \label{tab:comunication_semantics_profile}
\end{table}

This document provides the Interface Design Description IDD to the \textit{authorizationToken -- Service Description} document \cite{ca-authorizationToken-sd}.
The operations, their semantics and the data models are the same as in the HTTP/TLS/JSON IDD \cite{ca-authorizationToken-http-idd}; this document specifies the MQTT binding.

The MQTT interface is offered only when the system runs with \texttt{mqtt.api.enabled=true}; it is then registered with the interface template \texttt{generic\_\allowbreak{}mqtts} (\texttt{generic\_\allowbreak{}mqtt} without TLS), the broker address and port (\texttt{mqtt.\allowbreak{}broker.\allowbreak{}address}, \texttt{mqtt.\allowbreak{}broker.\allowbreak{}port}) and the base topic \texttt{arrowhead/\allowbreak{}consumer-\allowbreak{}authorization/\allowbreak{}authorization-\allowbreak{}token/\allowbreak{}}. In the Arrowhead-test-cloud-java deployment the MQTT interfaces of the core systems are disabled.

\newpage

\section{Service Operations}
\label{sec:operations}

The provider subscribes to one topic per operation, the base topic followed by the operation name:

\noindent\begin{longtable}{| p{5cm} | p{10.4cm} |} \hline
\rowcolor{gray!33} Operation & Request topic \\ \hline \endhead
\texttt{generate} & \texttt{arrowhead/\allowbreak{}consumer-\allowbreak{}authorization/\allowbreak{}authorization-\allowbreak{}token/\allowbreak{}generate} \\ \hline
\texttt{verify} & \texttt{arrowhead/\allowbreak{}consumer-\allowbreak{}authorization/\allowbreak{}authorization-\allowbreak{}token/\allowbreak{}verify} \\ \hline
\texttt{get-\allowbreak{}public-\allowbreak{}key} & \texttt{arrowhead/\allowbreak{}consumer-\allowbreak{}authorization/\allowbreak{}authorization-\allowbreak{}token/\allowbreak{}get-\allowbreak{}public-\allowbreak{}key} \\ \hline
\texttt{register-\allowbreak{}encryption-\allowbreak{}key} & \texttt{arrowhead/\allowbreak{}consumer-\allowbreak{}authorization/\allowbreak{}authorization-\allowbreak{}token/\allowbreak{}register-\allowbreak{}encryption-\allowbreak{}key} \\ \hline
\texttt{unregister-\allowbreak{}encryption-\allowbreak{}key} & \texttt{arrowhead/\allowbreak{}consumer-\allowbreak{}authorization/\allowbreak{}authorization-\allowbreak{}token/\allowbreak{}unregister-\allowbreak{}encryption-\allowbreak{}key} \\ \hline
\end{longtable}

A consumer publishes a request envelope (\mref{MqttRequestTemplate}) to the topic and receives the answer (\mref{MqttResponseTemplate}) on its own \texttt{responseTopic}. The \texttt{authentication} field carries the requester's identity according to the provider's authentication policy: \texttt{IDENTITY-TOKEN//<token>} (outsourced or internal policy), \texttt{SYSTEM//<name>} (declared policy) or the requester's X.509 certificate (certificate policy). The \texttt{status} of the response uses the HTTP status code numbers (200, 201, 204, 400, 401, 403, 404, 408, 423, 500, 503); on error the payload is an \mref{ErrorMessage}.

\begin{lstlisting}[caption={Request envelope.}]
{
  "traceId": "a1b2c3",
  "authentication": "IDENTITY-TOKEN//<token>",
  "responseTopic": "arrowhead/response/TemperatureDisplay/a1b2c3",
  "qosRequirement": 1,
  "params": {},
  "payload": {}
}
\end{lstlisting}

\begin{lstlisting}[caption={Response envelope.}]
{
  "status": 200,
  "traceId": "a1b2c3",
  "receiver": "TemperatureDisplay",
  "payload": {}
}
\end{lstlisting}

\osubsection{generate}{PUBLISH}{\texttt{arrowhead/\allowbreak{}consumer-\allowbreak{}authorization/\allowbreak{}authorization-\allowbreak{}token/\allowbreak{}generate}}{\mref[Authorization\allowbreak{}Token\allowbreak{}Generation\allowbreak{}Request]{AuthorizationTokenGenerationRequest}}{\mref[Authorization\allowbreak{}Token\allowbreak{}Generation\allowbreak{}Response]{AuthorizationTokenGenerationResponse}}

Generates a token of the requested variant for the requester (as consumer) to the provider's target and operation, if authorized.

\begin{lstlisting}[caption={Payload of an \oref{generate} request.}]
{
  "tokenVariant": "TIME_LIMITED_TOKEN_AUTH",
  "provider": "TemperatureSensor",
  "targetType": "SERVICE_DEF",
  "target": "temperatureReading",
  "scope": "read"
}
\end{lstlisting}

\osubsection{verify}{PUBLISH}{\texttt{arrowhead/\allowbreak{}consumer-\allowbreak{}authorization/\allowbreak{}authorization-\allowbreak{}token/\allowbreak{}verify}}{\pref{String}}{\mref[Authorization\allowbreak{}Token\allowbreak{}Verify\allowbreak{}Response]{AuthorizationTokenVerifyResponse}}

Verifies a token presented to the requester (as provider), counts a use of usage-limited tokens and returns its consumer, target and scope.

\begin{lstlisting}[caption={Payload of an \oref{verify} request.}]
"<token>"
\end{lstlisting}

\osubsection{get-public-key}{PUBLISH}{\texttt{arrowhead/\allowbreak{}consumer-\allowbreak{}authorization/\allowbreak{}authorization-\allowbreak{}token/\allowbreak{}get-\allowbreak{}public-\allowbreak{}key}}{}{\pref{String}}

Returns the public key (Base64 X.509 Subject\allowbreak{}Public\allowbreak{}Key\allowbreak{}Info) with which providers verify JSON Web Tokens.

\osubsection{register-encryption-key}{PUBLISH}{\texttt{arrowhead/\allowbreak{}consumer-\allowbreak{}authorization/\allowbreak{}authorization-\allowbreak{}token/\allowbreak{}register-\allowbreak{}encryption-\allowbreak{}key}}{\mref[Authorization\allowbreak{}Encryption\allowbreak{}Key\allowbreak{}Registration\allowbreak{}Request]{AuthorizationEncryptionKeyRegistrationRequest}}{\pref{String}}

Registers or replaces the requester's AES key for encrypting its self-contained tokens. Returns the initialization vector (CBC) or an empty string; 201 for a new key.

\begin{lstlisting}[caption={Payload of an \oref{register-encryption-key} request.}]
{
  "key": "0123456789abcdef",
  "algorithm": "AES/ECB/PKCS5Padding"
}
\end{lstlisting}

\osubsection{unregister-encryption-key}{PUBLISH}{\texttt{arrowhead/\allowbreak{}consumer-\allowbreak{}authorization/\allowbreak{}authorization-\allowbreak{}token/\allowbreak{}unregister-\allowbreak{}encryption-\allowbreak{}key}}{}{}

Removes the requester's encryption key. Returns 200 if removed, 204 if there was none.

\newpage

\section{Data Models}
\label{sec:model}

Here, all data objects that can be part of the service calls associated with this service are listed in alphabetic order.
Each \textit{struct} denotes a JSON \pref{Object} with the named fields; fields that are not needed may be omitted.

\msubsection{struct}{MqttRequestTemplate}

The request envelope.

\noindent\begin{tabularx}{\textwidth}{| p{4.25cm} | p{3.5cm} | X |} \hline
\rowcolor{gray!33} Object Field & Value Type & Description \\ \hline
"traceId" & \pref{String} & Identifier chosen by the requester, returned in the response. \\ \hline
"authentication" & \pref{String} & Identity of the requester (see Section \ref{sec:operations}). \\ \hline
"responseTopic" & \pref{String} & Topic on which the response is published. \\ \hline
"qosRequirement" & \pref{Number} & MQTT QoS of the response (0, 1 or 2). \\ \hline
"params" & \pref{Object}$<$\pref{String}$>$ & Optional parameters (the query parameters of the HTTP binding). \\ \hline
"payload" & \pref{Object} & The operation input. \\ \hline
\end{tabularx}

\msubsection{struct}{MqttResponseTemplate}

The response envelope.

\noindent\begin{tabularx}{\textwidth}{| p{4.25cm} | p{3.5cm} | X |} \hline
\rowcolor{gray!33} Object Field & Value Type & Description \\ \hline
"status" & \pref{Number} & Status code. \\ \hline
"traceId" & \pref{String} & The trace identifier of the request. \\ \hline
"receiver" & \pref{String} & Name of the requester. \\ \hline
"payload" & \pref{Object} & The operation output, or an \mref{ErrorMessage}. \\ \hline
\end{tabularx}

\msubsection{struct}{ErrorMessage}




\noindent\begin{tabularx}{\textwidth}{| >{\raggedright\arraybackslash}p{3.9cm} | >{\raggedright\arraybackslash\small}p{5.3cm} | X |} \hline
\rowcolor{gray!33} Object Field & Value Type & Description \\ \hline
"error\allowbreak{}Message" & \pref{String} &  \\ \hline
"error\allowbreak{}Code" & \pref{Number} &  \\ \hline
"exception\allowbreak{}Type" & \pref{ExceptionType} &  \\ \hline
"origin" & \pref{String} &  \\ \hline
\end{tabularx}



\msubsection{struct}{AuthorizationEncryptionKeyRegistrationRequest}

Registers the requester's key for encrypting its self-contained tokens.


\noindent\begin{tabularx}{\textwidth}{| >{\raggedright\arraybackslash}p{3.9cm} | >{\raggedright\arraybackslash\small}p{5.3cm} | X |} \hline
\rowcolor{gray!33} Object Field & Value Type & Description \\ \hline
"key" & \pref{String} & AES key of the provider (at least 16 bytes). \\ \hline
"algorithm" & \pref{String} & Encryption algorithm of the key, e.g. \texttt{AES/\allowbreak{}ECB/\allowbreak{}PKCS5Padding} or \texttt{AES/\allowbreak{}CBC/\allowbreak{}PKCS5Padding}. \\ \hline
\end{tabularx}



\msubsection{struct}{AuthorizationTokenGenerationRequest}

Request of the consumer for a token to one provider's service operation.


\noindent\begin{tabularx}{\textwidth}{| >{\raggedright\arraybackslash}p{3.9cm} | >{\raggedright\arraybackslash\small}p{5.3cm} | X |} \hline
\rowcolor{gray!33} Object Field & Value Type & Description \\ \hline
"token\allowbreak{}Variant" & \pref{String} & Requested token variant, as interface security policy (e.g. \texttt{TIME\_\allowbreak{}LIMITED\_\allowbreak{}TOKEN\_\allowbreak{}AUTH}, \texttt{RSA\_\allowbreak{}SHA512\_\allowbreak{}JSON\_\allowbreak{}WEB\_\allowbreak{}TOKEN\_\allowbreak{}AUTH}). \\ \hline
"provider" & \pref{String} & Provider system the token is for. \\ \hline
"target\allowbreak{}Type" & \pref{String} & Target type: \texttt{SERVICE\_\allowbreak{}DEF} or \texttt{EVENT\_\allowbreak{}TYPE}. \\ \hline
"target" & \pref{String} & Service definition (or event type). \\ \hline
"scope" & \pref{String} & Operation the token is for. \\ \hline
\end{tabularx}



\msubsection{struct}{AuthorizationTokenGenerationResponse}

A generated token.


\noindent\begin{tabularx}{\textwidth}{| >{\raggedright\arraybackslash}p{3.9cm} | >{\raggedright\arraybackslash\small}p{5.3cm} | X |} \hline
\rowcolor{gray!33} Object Field & Value Type & Description \\ \hline
"token\allowbreak{}Type" & \pref{AuthorizationTokenType} & Token type: \texttt{TIME\_\allowbreak{}LIMITED\_\allowbreak{}TOKEN}, \texttt{USAGE\_\allowbreak{}LIMITED\_\allowbreak{}TOKEN}, \texttt{SELF\_\allowbreak{}CONTAINED\_\allowbreak{}TOKEN} or \texttt{TRANSLATION\_\allowbreak{}BRIDGE\_\allowbreak{}TOKEN}. \\ \hline
"target\allowbreak{}Type" & \pref{AuthorizationTargetType} & Target type: \texttt{SERVICE\_\allowbreak{}DEF} or \texttt{EVENT\_\allowbreak{}TYPE}. \\ \hline
"token" & \pref{String} & The token (opaque string, Base64 or JSON Web Token). \\ \hline
"usage\allowbreak{}Limit" & \pref{Number} & Number of allowed uses (usage-limited tokens). \\ \hline
"expires\allowbreak{}At" & \pref{DateTime} & UTC expiry of time-limited tokens. \\ \hline
\end{tabularx}



\msubsection{struct}{AuthorizationTokenVerifyResponse}

Result of verifying a token presented to a provider.


\noindent\begin{tabularx}{\textwidth}{| >{\raggedright\arraybackslash}p{3.9cm} | >{\raggedright\arraybackslash\small}p{5.3cm} | X |} \hline
\rowcolor{gray!33} Object Field & Value Type & Description \\ \hline
"verified" & \pref{Boolean} & True when the token is valid. \\ \hline
"consumer\allowbreak{}Cloud" & \pref{String} & Cloud identifier of the consumer system. \\ \hline
"consumer" & \pref{String} & Name of the consumer system. \\ \hline
"target\allowbreak{}Type" & \pref{AuthorizationTargetType} & Target type: \texttt{SERVICE\_\allowbreak{}DEF} or \texttt{EVENT\_\allowbreak{}TYPE}. \\ \hline
"target" & \pref{String} & Target of the policy or token: a service definition or event type. \\ \hline
"scope" & \pref{String} & Operation the token or policy is limited to; empty for all operations. \\ \hline
\end{tabularx}



\subsection{Primitives}
\label{sec:model:primitives}

As all messages are encoded using the JSON format \cite{rfc8259}, the following primitive constructs, part of that standard, become available.

\noindent\begin{tabularx}{\textwidth}{| p{4.6cm} | X |} \hline
\rowcolor{gray!33} Type & Description \\ \hline
\pdef{AuthorizationTargetType} & One of \texttt{SERVICE\_\allowbreak{}DEF}, \texttt{EVENT\_\allowbreak{}TYPE} (JSON String). \\ \hline
\pdef{AuthorizationTokenType} & One of \texttt{TIME\_\allowbreak{}LIMITED\_\allowbreak{}TOKEN}, \texttt{USAGE\_\allowbreak{}LIMITED\_\allowbreak{}TOKEN}, \texttt{SELF\_\allowbreak{}CONTAINED\_\allowbreak{}TOKEN}, \texttt{TRANSLATION\_\allowbreak{}BRIDGE\_\allowbreak{}TOKEN} (JSON String). \\ \hline
\pdef{Boolean} & JSON Boolean, \texttt{true} or \texttt{false}. \\ \hline
\pdef{DateTime} & JSON String with an ISO 8601 UTC timestamp, e.g. \texttt{2026-10-22T12:00:00Z}. \\ \hline
\pdef{ExceptionType} & One of \texttt{ARROWHEAD}, \texttt{INVALID\_\allowbreak{}PARAMETER}, \texttt{AUTH}, \texttt{FORBIDDEN}, \texttt{DATA\_\allowbreak{}NOT\_\allowbreak{}FOUND}, \texttt{TIMEOUT}, \texttt{LOCKED}, \texttt{INTERNAL\_\allowbreak{}SERVER\_\allowbreak{}ERROR}, \texttt{EXTERNAL\_\allowbreak{}SERVER\_\allowbreak{}ERROR} (JSON String). \\ \hline
\pdef{Number} & JSON Number, an integer. \\ \hline
\pdef{Object} & JSON Object; \pref{Object}$<$A$>$ has string keys and values of type A. \\ \hline
\pdef{String} & JSON String. \\ \hline
\end{tabularx}

\newpage

\bibliographystyle{IEEEtran}
\bibliography{bibliography}

\newpage

\section{Revision History}
\subsection{Amendments}

\noindent\begin{tabularx}{\textwidth}{| p{1cm} | p{3cm} | p{2cm} | X | p{4cm} |} \hline
\rowcolor{gray!33} No. & Date & Version & Subject of Amendments & Author \\ \hline

1 & 2026-10-22 & \arrowversion & Document created from the Arrowhead LaTeX templates for the 5.2.2 Java Spring implementation, including the local extensions of this workspace. & Jerker Delsing \\ \hline
\end{tabularx}

\subsection{Quality Assurance}

\noindent\begin{tabularx}{\textwidth}{| p{1cm} | p{3cm} | p{2cm} | X |} \hline
\rowcolor{gray!33} No. & Date & Version & Approved by \\ \hline

1 & & \arrowversion  &  \\ \hline

\end{tabularx}

\end{document}
